% ===== Topic 2, part d: problems 41-50 (VI) =====
\providecommand{\tbGrid}[4]{\begin{tikzpicture}[baseline=0.2cm,x=0.5cm,y=0.5cm]
  \draw[Ink,line width=0.7pt] (0,0) grid[step=1] (2,2);
  \node[font=\sffamily\bfseries,text=tmBlue!80!black,inner sep=0pt,#4] at (#1,#2) {#3};\end{tikzpicture}}
% helper inside a tikzpicture: \rvbbG{x0}{y0}{cx}{cy}{letter}{transform} draws a 2x2 grid (cell 0.5) at (x0,y0)
% (all four outer edges drawn explicitly; `grid' can drop an edge through rounding)
\providecommand{\rvbbG}[6]{\draw[Ink,line width=0.7pt] (#1,#2) rectangle ++(1,1)
  ({#1+0.5},#2)--++(0,1) (#1,{#2+0.5})--++(1,0);
  \node[font=\sffamily\bfseries,text=tmBlue!80!black,inner sep=0pt,#6] at ({#1+#3},{#2+#4}) {#5};}

% ---------------------------------------------------------------- 41
\begin{probbox}{41}{\stars{1}}
\textit{(Mở rộng Bài 16.)} Ta chỉ xét các hình chữ nhật có độ dài hai cạnh là các \emph{số nguyên dương}
(đơn vị cm); hình vuông cũng được coi là hình chữ nhật.
\begin{enumerate}[label=\alph*),leftmargin=1.6em,itemsep=1pt]
\item Tìm tất cả các hình chữ nhật mà số đo diện tích (cm$^2$) \emph{bằng} số đo chu vi (cm).
\item Tìm tất cả các hình chữ nhật mà số đo diện tích \emph{gấp đôi} số đo chu vi.
\end{enumerate}
\end{probbox}
\probsol{
\textbf{Phân tích đề.} Gọi hai cạnh là $a$ và $b$ (số nguyên dương). Không mất tính tổng quát, giả sử
$a\le b$ ($a$ là cạnh ngắn). Diện tích là $ab$, chu vi là $2(a+b)$. Đây là bài toán ``tìm nghiệm nguyên'':
một phương trình có hai ẩn. Mẹo: dùng $a\le b$ để \emph{chặn} $a$ trong một khoảng nhỏ, rồi thử từng giá trị.

\textbf{a)} Điều kiện: $ab=2a+2b$.

\step{1}{Chặn trên cho $a$}
Vì $a\le b$ nên $2a+2b\le 2b+2b=4b$. Suy ra $ab\le 4b$. Chia hai vế cho $b>0$: $a\le 4$.

\step{2}{Chặn dưới cho $a$}
Vì $2a>0$ nên $ab=2a+2b>2b$. Chia hai vế cho $b>0$: $a>2$. Vậy $a=3$ hoặc $a=4$.

\step{3}{Thử từng giá trị}
\begin{itemize}[leftmargin=1.4em,itemsep=1pt]
\item $a=3$: $3b=6+2b\Rightarrow b=6$. Hình chữ nhật $3\times 6$.
\item $a=4$: $4b=8+2b\Rightarrow 2b=8\Rightarrow b=4$. Hình vuông $4\times 4$.
\end{itemize}
\begin{figC}\begin{scope}[x=0.45cm,y=0.45cm]
\fill[tmBlue!15] (0,0) rectangle (6,3);\draw[Line] (0,0) grid[step=1] (6,3);\draw[tmBlue!80!black,line width=1pt] (0,0) rectangle (6,3);
\node[below,font=\footnotesize] at (3,0){$S=18$, $C=18$};
\begin{scope}[xshift=4cm]
\fill[tmGreen!15] (0,0) rectangle (4,4);\draw[Line] (0,0) grid[step=1] (4,4);\draw[tmGreen!70!black,line width=1pt] (0,0) rectangle (4,4);
\node[below,font=\footnotesize] at (2,0){$S=16$, $C=16$};
\end{scope}
\end{scope}\end{figC}

\textbf{b)} Điều kiện: $ab=2\cdot 2(a+b)=4a+4b$.

\step{4}{Chặn $a$ như trên}
$ab=4a+4b\le 8b\Rightarrow a\le 8$; và $ab=4a+4b>4b\Rightarrow a>4$. Vậy $a\in\{5;6;7;8\}$.

\step{5}{Thử từng giá trị}
Từ $ab=4a+4b$ ta có $ab-4b=4a$, tức $b(a-4)=4a$, nên $b=\dfrac{4a}{a-4}$.
\begin{center}\begin{tabular}{|c|c|c|c|c|}\hline
$a$ & $5$ & $6$ & $7$ & $8$\\\hline
$b=\dfrac{4a}{a-4}$ & $\dfrac{20}{1}=20$ & $\dfrac{24}{2}=12$ & $\dfrac{28}{3}$ (không nguyên) & $\dfrac{32}{4}=8$\\\hline
\end{tabular}\end{center}
Có ba hình: $5\times 20$, $6\times 12$, $8\times 8$.

\begin{proofbox}Vì sao không sót? Mọi hình chữ nhật thỏa mãn đều có cạnh ngắn $a$ nằm trong khoảng đã chặn,
và với mỗi $a$, cạnh $b$ được xác định duy nhất từ phương trình. Ta đã thử hết các $a$ đó.\end{proofbox}
\begin{tipbox}Cách khác cho câu a): $ab-2a-2b+4=4$, tức là $(a-2)(b-2)=4=1\cdot 4=2\cdot 2$. Suy ra
$a-2=1,\ b-2=4$ hoặc $a-2=b-2=2$. Lại được $3\times 6$ và $4\times 4$.\end{tipbox}
\verify{$3\times 6$: $S=18$, $C=2\cdot 9=18$. $4\times4$: $16=16$. $5\times 20$: $S=100$, $C=50$;
$6\times 12$: $S=72$, $C=36$; $8\times 8$: $S=64$, $C=32$ -- đều gấp đôi.}
\ans{a) $3\times 6$ và $4\times 4$. \quad b) $5\times 20$, $6\times 12$, $8\times 8$.}
}

% ---------------------------------------------------------------- 42
\begin{probbox}{42}{\stars{1}}
\textit{(Mở rộng Bài 17.)} Chữ P được viết ở ô trên -- phải của lưới $2\times 2$: \tbGrid{1.5}{1.5}{P}{}\,.
Ta được dùng ba phép biến hình, bao nhiêu lần và theo thứ tự nào cũng được:
\begin{itemize}[leftmargin=1.4em,itemsep=1pt]
\item $Q$: quay lưới $90^\circ$ theo chiều kim đồng hồ quanh tâm;
\item $D$: lật lưới qua đường thẳng \emph{đứng} đi qua tâm;
\item $N$: lật lưới qua đường thẳng \emph{nằm ngang} đi qua tâm.
\end{itemize}
\begin{enumerate}[label=\alph*),leftmargin=1.6em,itemsep=1pt]
\item Vẽ kết quả khi thực hiện $D$ rồi $N$. Có thể nhận được kết quả đó chỉ bằng các phép $Q$ không?
\item Chứng tỏ rằng thực hiện $D$ rồi $Q$ cho cùng kết quả với một phép lật qua một \emph{đường chéo} của lưới.
Đó là đường chéo nào?
\item Có tất cả bao nhiêu hình khác nhau có thể nhận được (kể cả hình ban đầu)? Hãy giải thích vì sao không
thể nhiều hơn.
\end{enumerate}
\end{probbox}
\probsol{
\textbf{Phân tích đề.} Như trong Bài 17, ta coi lưới là một \emph{tấm thẻ} vuông trong suốt, có vẽ chữ P. Mỗi
phép biến hình là một cách nhấc tấm thẻ lên rồi đặt lại vừa khít chỗ cũ. Để theo dõi, ta ghi nhớ hai điều:
chữ P nằm ở ô nào, và nó bị quay/lật ra sao.

\textbf{a)}
\step{1}{Thực hiện $D$}
Lật qua đường đứng: ô trên -- phải sang ô trên -- trái, chữ P bị ``soi gương'' (bụng quay sang trái).
\step{2}{Thực hiện tiếp $N$}
Lật qua đường ngang: ô trên -- trái xuống ô dưới -- trái, chữ bị lật ngược trên -- dưới. Kết quả: chữ P ở ô
dưới -- trái, \emph{lộn đầu} (giống chữ ``d''), như hình thứ ba bên dưới.
\begin{figC}
\rvbbG{0}{0}{0.75}{0.75}{P}{}\node at (1.75,0.5){$\xrightarrow{\;D\;}$};
\rvbbG{2.5}{0}{0.25}{0.75}{P}{xscale=-1}\node at (4.25,0.5){$\xrightarrow{\;N\;}$};
\rvbbG{5}{0}{0.25}{0.25}{P}{rotate=180}
\node[anchor=west] at (6.2,0.5){\small$=$ quay $180^\circ$};
\end{figC}
\step{3}{So sánh với các phép $Q$}
Quay $Q$ hai lần (tức quay $180^\circ$): ô trên -- phải đi tới ô đối diện qua tâm là ô dưới -- trái, và chữ bị
lộn đầu. Trùng với kết quả trên. Vậy \textbf{$D$ rồi $N$ cho cùng kết quả với $Q$ rồi $Q$.}
\begin{recallbox}Hai lần lật qua hai đường vuông góc với nhau (cùng đi qua tâm) cho kết quả bằng một phép quay
$180^\circ$ quanh tâm.\end{recallbox}

\textbf{b)}
\step{4}{Thực hiện $D$ rồi $Q$}
$D$: chữ P sang ô trên -- trái, bị soi gương. $Q$ (quay $90^\circ$ theo chiều kim đồng hồ): ô trên -- trái sang ô
trên -- phải. Kết quả: P \emph{trở về ô trên -- phải} nhưng thân nằm ngang, bụng ở phía trên:
\begin{figC}
\rvbbG{0}{0}{0.75}{0.75}{P}{}\node at (1.75,0.5){$\xrightarrow{\;D\;}$};
\rvbbG{2.5}{0}{0.25}{0.75}{P}{xscale=-1}\node at (4.25,0.5){$\xrightarrow{\;Q\;}$};
\draw[tmOrange,dashed,line width=0.9pt] (4.8,-0.2)--(6.2,1.2);
\rvbbG{5}{0}{0.75}{0.75}{P}{cm={0,1,1,0,(0,0)},fill=white,inner sep=0.8pt}
\end{figC}
\step{5}{Nhận ra phép lật qua đường chéo}
Ô trên -- phải không đổi chỗ, nghĩa là nó nằm trên trục lật. Đường chéo đi qua ô trên -- phải (và ô dưới -- trái)
là đường chéo nối \textbf{góc dưới -- trái với góc trên -- phải} (nét đứt). Lật qua đường này: hướng ``lên''
thành hướng ``sang phải'' và ngược lại, nên thân P (thẳng đứng) thành nằm ngang, bụng P (bên phải) thành ở
phía trên -- đúng như hình. Có thể kiểm tra thêm: một chữ đặt ở ô trên -- trái sẽ chuyển sang ô dưới -- phải
trong cả hai cách.

\textbf{c)}
\step{6}{Không thể có quá $8$ hình}
\begin{proofbox}Đánh dấu góc trên -- phải của tấm thẻ bằng một ngôi sao. Mỗi lần đặt lại, ngôi sao chỉ có thể
nằm ở $1$ trong $4$ góc của khung. Ngoài ra tấm thẻ chỉ có thể đang \emph{ngửa} hoặc \emph{úp} (mỗi lần lật
$D$ hay $N$ đổi ngửa thành úp và ngược lại; quay $Q$ không đổi). Khi đã biết vị trí ngôi sao và thẻ ngửa hay
úp thì vị trí của tấm thẻ, và do đó cả hình, được xác định hoàn toàn. Vậy có nhiều nhất $4\times 2=8$
hình.\end{proofbox}
\step{7}{Có đủ $8$ hình}
Bốn hình ``thẻ ngửa'': quay $0^\circ$, $90^\circ$, $180^\circ$, $270^\circ$ (dùng $Q$ từ $0$ đến $3$ lần).
Bốn hình ``thẻ úp'': mỗi hình trên thực hiện thêm $D$. Chữ P không có trục đối xứng nào, nên $8$ hình này
khác nhau từng đôi một:
\begin{figC}
\rvbbG{0}{0}{0.75}{0.75}{P}{}\rvbbG{1.5}{0}{0.75}{0.25}{P}{rotate=-90}
\rvbbG{3}{0}{0.25}{0.25}{P}{rotate=180}\rvbbG{4.5}{0}{0.25}{0.75}{P}{rotate=90}
\rvbbG{0}{-1.5}{0.25}{0.75}{P}{xscale=-1}\rvbbG{1.5}{-1.5}{0.25}{0.25}{P}{cm={0,-1,-1,0,(0,0)}}
\rvbbG{3}{-1.5}{0.75}{0.25}{P}{yscale=-1}\rvbbG{4.5}{-1.5}{0.75}{0.75}{P}{cm={0,1,1,0,(0,0)}}
\node[anchor=east,font=\footnotesize] at (-0.2,0.5){thẻ ngửa};\node[anchor=east,font=\footnotesize] at (-0.2,-1){thẻ úp};
\end{figC}
\begin{warnbox}Đừng đếm ``$4$ ô $\times$ $8$ tư thế $=32$''. Vị trí ô và tư thế của chữ \emph{gắn liền} với
nhau, vì chúng đi theo cùng một tấm thẻ.\end{warnbox}
\ans{a) P ở ô dưới -- trái, lộn đầu; bằng $Q$ rồi $Q$ (quay $180^\circ$). \ b) Đường chéo nối góc dưới -- trái
với góc trên -- phải. \ c) $8$ hình.}
}

% ---------------------------------------------------------------- 43
\begin{probbox}{43}{\stars{1}}
\textit{(Mở rộng Bài 18.)} Cho hình chữ nhật $ABCD$ với $AB=12$ cm, $BC=8$ cm. Điểm $P$ nằm bên trong hình chữ nhật.
\begin{figC}\begin{scope}[x=0.36cm,y=0.36cm]
\fill[tmBlue!25] (0,0)--(12,0)--(7,3)--cycle;\fill[tmOrange!25] (12,8)--(0,8)--(7,3)--cycle;
\draw[Ink,line width=1pt] (0,0) rectangle (12,8);
\draw[Ink] (0,0)--(7,3)--(12,0) (0,8)--(7,3)--(12,8);
\fill[Ink] (7,3) circle (2.5pt);\node[right=2pt] at (7,3){$P$};
\node[below left] at (0,0){$A$};\node[below right] at (12,0){$B$};\node[above right] at (12,8){$C$};\node[above left] at (0,8){$D$};
\end{scope}\end{figC}
\begin{enumerate}[label=\alph*),leftmargin=1.6em,itemsep=1pt]
\item Chứng minh rằng dù $P$ ở đâu, ta luôn có $S_{PAB}+S_{PCD}=48$ cm$^2$.
\item Biết $S_{PAB}=18$ cm$^2$ và $S_{PBC}=20$ cm$^2$. Tính $S_{PCD}$, $S_{PDA}$ và khoảng cách từ $P$ đến
bốn cạnh của hình chữ nhật.
\item Có điểm $P$ nào bên trong hình chữ nhật mà $S_{PAB}=20$, $S_{PBC}=30$ và $S_{PCD}=30$ (cm$^2$) không?
\end{enumerate}
\end{probbox}
\probsol{
\textbf{Phân tích đề.} Cả bốn tam giác đều có đáy là một cạnh hình chữ nhật và đỉnh là $P$. Chiều cao của
mỗi tam giác chính là \emph{khoảng cách từ $P$ đến cạnh đó}. Ta gọi khoảng cách từ $P$ đến $AB$, $BC$, $CD$,
$DA$ lần lượt là $h_1,h_2,h_3,h_4$.
\begin{figC}\begin{scope}[x=0.36cm,y=0.36cm]
\draw[Ink,line width=1pt] (0,0) rectangle (12,8);
\draw[tmBlue,thick,<->] (7,0.1)--(7,2.9);\draw[tmOrange!80!black,thick,<->] (7,3.1)--(7,7.9);
\draw[tmGreen!70!black,thick,<->] (7.1,3)--(11.9,3);\draw[tmPurple,thick,<->] (0.1,3)--(6.9,3);
\fill[Ink] (7,3) circle (2.5pt);
\node[right,font=\footnotesize] at (7,1.5){$h_1$};\node[right,font=\footnotesize] at (7,5.5){$h_3$};
\node[above,font=\footnotesize] at (9.5,3){$h_2$};\node[above,font=\footnotesize] at (3.5,3){$h_4$};
\node[below left] at (0,0){$A$};\node[below right] at (12,0){$B$};\node[above right] at (12,8){$C$};\node[above left] at (0,8){$D$};
\end{scope}\end{figC}

\textbf{a)}
\step{1}{Tổng hai chiều cao}
Kẻ qua $P$ đường thẳng vuông góc với $AB$; vì $AB\parallel CD$ nên nó cũng vuông góc với $CD$. Đoạn nằm giữa
$AB$ và $CD$ có độ dài bằng $BC=8$, và $P$ chia đoạn này thành hai phần $h_1$ và $h_3$. Vậy $h_1+h_3=8$.
Tương tự $h_2+h_4=AB=12$.
\step{2}{Cộng hai diện tích}
\begin{recallbox}Diện tích tam giác $=\tfrac12\times$ đáy $\times$ chiều cao.\end{recallbox}
\[ S_{PAB}+S_{PCD}=\frac{12\cdot h_1}{2}+\frac{12\cdot h_3}{2}=6(h_1+h_3)=6\cdot 8=48 \text{ (cm}^2). \]
Kết quả không phụ thuộc vào vị trí của $P$. (Đó đúng bằng nửa diện tích hình chữ nhật $12\cdot 8=96$.)
Tương tự, $S_{PBC}+S_{PDA}=\tfrac{8}{2}(h_2+h_4)=4\cdot 12=48$ cm$^2$.

\textbf{b)}
\step{3}{Dùng kết quả câu a}
$S_{PCD}=48-18=30$ cm$^2$, \quad $S_{PDA}=48-20=28$ cm$^2$.
\step{4}{Tính bốn khoảng cách}
Từ $S=\tfrac12\cdot$ đáy $\cdot h$ suy ra $h=2S:$ đáy.
\[ h_1=\frac{2\cdot 18}{12}=3,\quad h_2=\frac{2\cdot 20}{8}=5,\quad h_3=\frac{2\cdot 30}{12}=5,\quad h_4=\frac{2\cdot 28}{8}=7 \text{ (cm)}. \]

\textbf{c)}
\step{5}{Kiểm tra điều kiện bắt buộc}
Theo câu a), luôn có $S_{PAB}+S_{PCD}=48$. Nhưng $20+30=50\ne 48$. Vậy \textbf{không có} điểm $P$ nào như thế.

\begin{tipbox}Bài 18 là trường hợp đặc biệt: $P$ là tâm hình vuông nên bốn tam giác bằng nhau. Ở đây $P$ bất kì
nhưng ``tổng hai tam giác đối diện'' vẫn luôn bằng nửa hình.\end{tipbox}
\verify{$h_1+h_3=3+5=8=BC$ và $h_2+h_4=5+7=12=AB$. Đặt $A(0;0)$ thì $P(7;3)$ -- đúng là điểm trong hình vẽ.}
\ans{b) $S_{PCD}=30$ cm$^2$, $S_{PDA}=28$ cm$^2$; khoảng cách đến $AB,BC,CD,DA$ là $3;5;5;7$ cm. \ c) Không có.}
}

% ---------------------------------------------------------------- 44
\begin{probbox}{44}{\stars{1}}
\textit{(Mở rộng Bài 22.)} Hình bình hành $ABCD$ có diện tích $48$ cm$^2$.
\begin{enumerate}[label=\alph*),leftmargin=1.6em,itemsep=1pt]
\item Gọi $M$, $N$ lần lượt là trung điểm của $BC$ và $CD$. Tính diện tích tam giác $AMN$.
\item Giữ $N$ là trung điểm của $CD$, nhưng lấy $M$ trên cạnh $BC$ sao cho $BM=3\,MC$. Tính diện tích tam giác $AMN$.
\end{enumerate}
\begin{figC}\begin{scope}[x=0.55cm,y=0.55cm]
\fill[tmBlue!25] (0,0)--(7,2)--(5,4)--cycle;
\draw[Ink,line width=1pt] (0,0)--(6,0)--(8,4)--(2,4)--cycle;
\draw[tmBlue!80!black,line width=1pt] (0,0)--(7,2)--(5,4)--cycle;
\node[below left] at (0,0){$A$};\node[below right] at (6,0){$B$};\node[above right] at (8,4){$C$};\node[above left] at (2,4){$D$};
\node[right] at (7,2){$M$};\node[above] at (5,4){$N$};
\end{scope}\end{figC}
\end{probbox}
\probsol{
\textbf{Phân tích đề.} Tam giác $AMN$ nằm ``chéo'' bên trong hình bình hành. Ta lấy cả hình bình hành bỏ đi ba
tam giác ở các góc: $ABM$, $MCN$ và $NDA$.
\[ S_{AMN}=S_{ABCD}-S_{ABM}-S_{MCN}-S_{NDA}. \]
\begin{recallbox}(1) Mỗi đường chéo chia hình bình hành thành hai tam giác có diện tích bằng nhau ($=24$ cm$^2$).
(2) Hai tam giác chung đỉnh, có hai đáy nằm trên cùng một đường thẳng thì tỉ số diện tích bằng tỉ số hai đáy.
\end{recallbox}

\textbf{a)}
\step{1}{Tam giác $ABM$}
So với tam giác $ABC$ ($=24$): chung đỉnh $A$, đáy $BM=\tfrac12 BC$. Vậy $S_{ABM}=\tfrac12\cdot 24=12$.
\step{2}{Tam giác $NDA$}
So với tam giác $ACD$ ($=24$): chung đỉnh $A$, đáy $DN=\tfrac12 DC$. Vậy $S_{NDA}=12$.
\step{3}{Tam giác $MCN$ (làm hai bước, như Bài 22)}
Tam giác $BCD$ có diện tích $24$. Tam giác $MCD$ chung đỉnh $D$ với nó, đáy $CM=\tfrac12 CB$: $S_{MCD}=12$.
Tam giác $MCN$ chung đỉnh $M$ với tam giác $MCD$, đáy $CN=\tfrac12 CD$: $S_{MCN}=\tfrac12\cdot 12=6$.
\step{4}{Kết quả}
$S_{AMN}=48-12-6-12=18$ cm$^2$.

\textbf{b)}
\step{5}{Làm lại với $MC=\tfrac14 BC$}
Vì $BM=3\,MC$ nên $BC=4\,MC$, tức $MC=\tfrac14 BC$ và $BM=\tfrac34 BC$.
\begin{itemize}[leftmargin=1.4em,itemsep=1pt]
\item $S_{ABM}=\tfrac34\cdot 24=18$;
\item $S_{NDA}=12$ (không đổi);
\item $S_{MCD}=\tfrac14\cdot 24=6$, rồi $S_{MCN}=\tfrac12\cdot 6=3$.
\end{itemize}
$S_{AMN}=48-18-3-12=15$ cm$^2$.

\begin{tipbox}Quy tắc ``góc'': nếu trên hai cạnh của một tam giác, tính từ cùng một đỉnh, ta lấy các phần
bằng $\tfrac{p}{q}$ và $\tfrac{r}{s}$ của cạnh, thì tam giác nhỏ ở góc đó có diện tích bằng
$\tfrac{p}{q}\cdot\tfrac{r}{s}$ tam giác lớn. Ở câu b): $\tfrac14\cdot\tfrac12=\tfrac18$ của $24$ là $3$.\end{tipbox}
\verify{Đặt $A(0;0)$, $B(8;0)$, $D(2;6)$, $C(10;6)$ (diện tích $8\cdot 6=48$). Câu a): $M(9;3)$, $N(6;6)$; bao tam giác
$AMN$ bằng hình chữ nhật $9\times 6=54$, trừ $\tfrac{9\cdot3}{2}=13{,}5$, $\tfrac{3\cdot 3}{2}=4{,}5$,
$\tfrac{6\cdot6}{2}=18$: còn $18$. Đúng.}
\ans{a) $18$ cm$^2$. \quad b) $15$ cm$^2$.}
}

% ---------------------------------------------------------------- 45
\begin{probbox}{45}{\stars{1}}
\textit{(Mở rộng Bài 21.)} Xét các tam giác có độ dài ba cạnh là các số nguyên dương và chu vi bằng $15$
(hai tam giác có cùng bộ ba độ dài cạnh chỉ tính là một).
\begin{enumerate}[label=\alph*),leftmargin=1.6em,itemsep=1pt]
\item Gọi $c$ là cạnh dài nhất. Chứng minh rằng $5\le c\le 7$.
\item Có tất cả bao nhiêu tam giác như vậy? Hãy liệt kê.
\item Trong đó có bao nhiêu tam giác cân? Bao nhiêu tam giác có \emph{đúng} hai cạnh bằng nhau?
\end{enumerate}
\end{probbox}
\probsol{
\textbf{Phân tích đề.} Gọi ba cạnh là $a\le b\le c$ (sắp xếp từ nhỏ đến lớn để không đếm trùng). Ta có
$a+b+c=15$, và theo bất đẳng thức tam giác chỉ cần kiểm tra $a+b>c$.
\begin{recallbox}Ba đoạn $a\le b\le c$ là ba cạnh của một tam giác khi và chỉ khi $a+b>c$ (tổng hai cạnh nhỏ
lớn hơn cạnh lớn nhất).\end{recallbox}

\textbf{a)}
\step{1}{Chặn trên}
$a+b=15-c$. Điều kiện $a+b>c$ trở thành $15-c>c$, tức $2c<15$, $c<7{,}5$. Vì $c$ nguyên nên $c\le 7$.
\step{2}{Chặn dưới}
$c$ là cạnh lớn nhất nên $a\le c$, $b\le c$, suy ra $15=a+b+c\le 3c$, tức $c\ge 5$.

\textbf{b)}
\step{3}{Liệt kê theo $c$}
Với mỗi $c$, ta có $a+b=15-c$ và $a\le b\le c$.
\begin{center}\begin{tabular}{|c|c|l|c|}\hline
$c$ & $a+b$ & Các cặp $(a;b)$ với $a\le b\le c$ & Số tam giác\\\hline
$7$ & $8$ & $(1;7),\ (2;6),\ (3;5),\ (4;4)$ & $4$\\
$6$ & $9$ & $(3;6),\ (4;5)$ \quad (không lấy $(2;7)$ vì $7>6$) & $2$\\
$5$ & $10$ & $(5;5)$ & $1$\\\hline
\end{tabular}\end{center}
Tổng cộng $4+2+1=7$ tam giác:
\[ (1;7;7),\ (2;6;7),\ (3;5;7),\ (4;4;7),\ (3;6;6),\ (4;5;6),\ (5;5;5). \]
\begin{warnbox}Với $c=7$, cặp $(1;7)$ cho tam giác $1;7;7$ -- rất ``gầy'' nhưng vẫn hợp lệ vì $1+7=8>7$.
Ngược lại, không có tam giác nào có cạnh $8$ vì khi đó hai cạnh còn lại có tổng $7<8$.\end{warnbox}

\textbf{c)}
\step{4}{Chọn ra các tam giác cân}
Tam giác cân có ít nhất hai cạnh bằng nhau: $(1;7;7)$, $(4;4;7)$, $(3;6;6)$, $(5;5;5)$ -- có $4$ tam giác.
Trong đó $(5;5;5)$ là tam giác đều (ba cạnh bằng nhau), nên có đúng hai cạnh bằng nhau là $3$ tam giác.
\verify{Mỗi bộ đều có tổng $15$ và thỏa $a+b>c$: $8>7$, $8>7$, $8>7$, $8>7$, $9>6$, $9>6$, $10>5$. Đúng.}
\ans{a) $5\le c\le 7$. \ b) $7$ tam giác. \ c) $4$ tam giác cân, trong đó $3$ tam giác có đúng hai cạnh bằng nhau.}
}

% ---------------------------------------------------------------- 46
\begin{probbox}{46}{\stars{2}}
\textit{(Mở rộng Bài 23.)} Trên lưới ô vuông (mỗi ô cạnh $1$), cho tam giác $ABC$ với $A(0;0)$, $B(12;0)$,
$C(3;9)$. Có đúng $43$ giao điểm của các đường kẻ lưới nằm \emph{bên trong} (không nằm trên cạnh) tam giác.
Với mỗi giao điểm $P$ như vậy, ta được ba tam giác $PAB$, $PBC$, $PCA$.
\begin{figC}\begin{scope}[x=0.36cm,y=0.36cm]
\draw[Line,line width=0.5pt] (0,0) grid[step=1] (12,9);
\draw[Ink,line width=1.3pt] (0,0)--(12,0)--(3,9)--cycle;
\node[below left] at (0,0){$A$};\node[below right] at (12,0){$B$};\node[above] at (3,9){$C$};
\end{scope}\end{figC}
\begin{enumerate}[label=\alph*),leftmargin=1.6em,itemsep=1pt]
\item Tính diện tích tam giác $ABC$.
\item Tìm điểm $P$ sao cho ba tam giác $PAB$, $PBC$, $PCA$ có diện tích bằng nhau.
\item Trong tất cả $3\cdot 43=129$ tam giác, có bao nhiêu tam giác có diện tích đúng bằng $\tfrac13$ diện tích
tam giác $ABC$? Có bao nhiêu vị trí của $P$ cho ít nhất một tam giác như thế?
\end{enumerate}
\end{probbox}
\probsol{
\textbf{Phân tích đề.} Đây là ``phiên bản một phần ba'' của Bài 23. Ý chính vẫn là: tam giác có đáy cố định thì
diện tích chỉ phụ thuộc vào chiều cao, và các điểm có cùng chiều cao nằm trên một đường thẳng song song với đáy.

\textbf{a)}
\step{1}{Diện tích tam giác $ABC$}
Đáy $AB=12$ nằm trên trục hoành, chiều cao là tung độ của $C$, bằng $9$:
$S_{ABC}=\tfrac{12\cdot 9}{2}=54$. Một phần ba là $18$.

\textbf{c)} (làm trước, câu b) suy ra ngay)
\step{2}{Tam giác $PAB$ có diện tích $18$}
Đáy $AB=12$ nên chiều cao phải là $2\cdot 18:12=3$, tức $P$ nằm trên đường ngang $y=3$.
Trên đường này: cạnh $AC$ (đi ``sang phải $1$, lên $3$'') qua điểm $(1;3)$; cạnh $BC$ (đi ``sang trái $1$,
lên $1$'' từ $B$) qua điểm $(9;3)$. Các điểm bên trong: $x=2,3,\dots,8$, tức $7$ điểm.

\step{3}{Tam giác $PBC$ có diện tích $18$}
\begin{proofbox}Lấy $E$ trên $AB$ với $EB=\tfrac13 AB=4$, tức $E(8;0)$. Tam giác $EBC$ và $ABC$ chung đỉnh $C$,
đáy $EB=\tfrac13 AB$ nên $S_{EBC}=\tfrac13\cdot 54=18$. Do đó $S_{PBC}=18$ khi và chỉ khi $P$ cách $BC$ đúng
bằng $E$, tức $P$ nằm trên đường thẳng qua $E$ song song với $BC$ (về phía trong tam giác).\end{proofbox}
$BC$ có hướng ``sang trái $1$, lên $1$''. Từ $E(8;0)$ đi theo hướng đó: $(7;1)$, $(6;2)$, $(5;3)$, $(4;4)$, $(3;5)$,
rồi $(2;6)$ -- điểm này nằm trên $AC$ (vì trên $AC$, tung độ gấp $3$ lần hoành độ), không tính. Có $5$ điểm.

\step{4}{Tam giác $PCA$ có diện tích $18$}
Tương tự, lấy $F$ trên $AB$ với $AF=4$, tức $F(4;0)$; $P$ phải nằm trên đường thẳng qua $F$ song song với $AC$
(hướng ``sang phải $1$, lên $3$''): $(5;3)$, rồi $(6;6)$ -- nằm trên $BC$ (vì $6+6=12$), không tính. Có $1$ điểm.
\begin{figC}\begin{scope}[x=0.36cm,y=0.36cm]
\draw[Line,line width=0.5pt] (0,0) grid[step=1] (12,9);
\draw[Ink,line width=1.3pt] (0,0)--(12,0)--(3,9)--cycle;
\draw[tmBlue,line width=1pt] (1,3)--(9,3);
\draw[tmOrange!90!black,line width=1pt] (8,0)--(2,6);
\draw[tmGreen!70!black,line width=1pt] (4,0)--(6,6);
\foreach \p in {(2,3),(3,3),(4,3),(6,3),(7,3),(8,3)}{\fill[tmBlue] \p circle (3pt);}
\foreach \p in {(7,1),(6,2),(4,4),(3,5)}{\fill[tmOrange!90!black] \p circle (3pt);}
\fill[AnsRed] (5,3) circle (4pt);
\node[below left] at (0,0){$A$};\node[below right] at (12,0){$B$};\node[above] at (3,9){$C$};
\node[below] at (8,0){$E$};\node[below] at (4,0){$F$};
\end{scope}\end{figC}

\step{5}{Đếm}
Số tam giác: $7+5+1=13$. Điểm $(5;3)$ nằm trên cả ba đường (điểm đỏ), nên được đếm ba lần. Số vị trí khác
nhau của $P$: $13-2=11$.

\textbf{b)}
\step{6}{Ba tam giác bằng nhau}
Ba tam giác có tổng diện tích bằng $54$, nên bằng nhau khi mỗi tam giác bằng $18$, tức $P$ phải nằm trên
\emph{cả ba} đường thẳng vừa tìm. Chỉ có một điểm như vậy: $P(5;3)$.

\begin{tipbox}Điểm $(5;3)$ chính là \emph{trọng tâm} của tam giác: hoành độ $\tfrac{0+12+3}{3}=5$, tung độ
$\tfrac{0+0+9}{3}=3$. Các em sẽ học kĩ về trọng tâm trong chương tam giác.\end{tipbox}
\verify{Với $P(5;3)$: $S_{PAB}=\tfrac{12\cdot 3}{2}=18$. $S_{PCA}$: bao bằng hình chữ nhật $5\times 9=45$, trừ
$\tfrac{3\cdot 9}{2}=13{,}5$ (góc trên -- trái), $\tfrac{2\cdot 6}{2}=6$ (góc trên -- phải),
$\tfrac{5\cdot 3}{2}=7{,}5$ (góc dưới -- phải): còn $18$. Vậy $S_{PBC}=54-18-18=18$. Đúng.}
\ans{a) $54$. \ b) $P(5;3)$. \ c) $13$ tam giác, ứng với $11$ vị trí của $P$.}
}

% ---------------------------------------------------------------- 47
\begin{probbox}{47}{\stars{1}}
\textit{(Mở rộng Bài 25.)} Ta lại xếp các hình chữ nhật $4\times 2$ thành ``cầu thang'' như Bài 25, nhưng lần này
đỉnh trên -- trái của mỗi hình (từ hình thứ hai) đặt tại điểm trên cạnh dưới của hình ngay phía trên, cách đỉnh
dưới -- trái của hình đó một khoảng $s$, với $s$ là $1$, $2$ hoặc $3$. (Bài 25 ứng với $s=2$.)
\begin{figC}\begin{scope}[x=0.38cm,y=0.38cm]
\foreach \i in {0,1,2,3}{\draw[Ink,line width=0.9pt] (\i,-2*\i) rectangle (\i+4,-2*\i+2);}
\draw[tmBlue,<->] (0,-0.5)--(1,-0.5);\node[below,font=\footnotesize] at (0.5,-0.5){$s$};
\node[above] at (2,2){$4$};\node[left] at (0,1){$2$};
\foreach \k in {0,1,2}{\fill[Ink] (5.5+0.5*\k,-6.6-0.5*\k) circle (2.3pt);}
\end{scope}\end{figC}
\begin{enumerate}[label=\alph*),leftmargin=1.6em,itemsep=1pt]
\item Với $s=1$ và $10$ hình chữ nhật, tính chu vi của cả hình.
\item Tìm công thức tính chu vi theo số hình $n$ và khoảng $s$.
\item Có thể chọn $s\in\{1;2;3\}$ và số hình $n$ để chu vi đúng bằng $100$ không? Tìm tất cả các cách.
\end{enumerate}
\end{probbox}
\probsol{
\textbf{Phân tích đề.} Ở Bài 25 ta đã thấy chu vi ``cầu thang'' bằng chu vi hình chữ nhật nhỏ nhất bao quanh
nó. Lập luận đó vẫn đúng với mọi $s$ (vì đường viền trên -- phải chỉ đi sang phải và đi xuống; đường viền dưới --
trái cũng vậy). Vậy chỉ cần tính bề ngang và chiều cao của cả hình.
\begin{recallbox}Hình ``cầu thang'' (đường viền chỉ đi sang phải/xuống ở mỗi phía) có chu vi bằng chu vi hình chữ
nhật bao quanh nó: $2\times(\text{bề ngang}+\text{chiều cao})$.\end{recallbox}

\textbf{b)} (làm trước vì câu a) suy ra ngay)
\step{1}{Bề ngang và chiều cao}
Mỗi hình mới lùi sang phải $s$. Hình đầu chiếm bề ngang $4$, mỗi hình sau thêm $s$:
bề ngang $=4+(n-1)s$. Mỗi hình mới thấp xuống $2$: chiều cao $=2n$.
\step{2}{Công thức}
\[ C=2\big(4+(n-1)s+2n\big). \]

\textbf{a)}
\step{3}{Thay $s=1$, $n=10$}
Bề ngang $4+9\cdot 1=13$, chiều cao $20$; $C=2(13+20)=66$.

\textbf{c)}
\step{4}{Xét từng giá trị của $s$}
\begin{itemize}[leftmargin=1.4em,itemsep=1pt]
\item $s=1$: $C=2(4+n-1+2n)=2(3n+3)=6n+6$. $6n+6=100\Rightarrow 6n=94$: không chia hết, loại.
\item $s=2$: $C=2(4+2n-2+2n)=2(4n+2)=8n+4$. $8n+4=100\Rightarrow n=12$. Nhận.
\item $s=3$: $C=2(4+3n-3+2n)=2(5n+1)=10n+2$. $10n+2=100\Rightarrow 10n=98$: loại.
\end{itemize}
Chỉ có một cách: $s=2$ và $n=12$.
\begin{warnbox}Nhớ rằng chiều cao $2n$ không phụ thuộc $s$; chỉ bề ngang thay đổi. Nhiều bạn cộng nhầm cả
những đoạn bị che bên trong.\end{warnbox}
\verify{Câu a) đếm trực tiếp: cạnh đứng $10\cdot 4=40$; cạnh ngang phía trên $4+9\cdot 1=13$; phía dưới
$9\cdot 1+4=13$. Tổng $66$. Đúng. Câu c): $s=2$, $n=12$: $8\cdot 12+4=100$. Đúng.}
\ans{a) $66$. \ b) $C=2\big(4+(n-1)s+2n\big)$. \ c) Chỉ có $s=2$, $n=12$.}
}

% ---------------------------------------------------------------- 48
\begin{probbox}{48}{\stars{1}}
\textit{(Mở rộng Bài 27.)} Hình dưới đây gồm $15$ hình bình hành nhỏ giống nhau, xếp thành $3$ hàng, mỗi hàng
$5$ hình. Hình bình hành nhỏ ở chính giữa được tô màu.
\begin{figC}\begin{scope}[x=0.5cm,y=0.5cm]
\fill[tmOrange!45] (5.4,1)--(7.6,1)--(8.6,2)--(6.4,2)--cycle;
\foreach \k in {0,...,5}{\draw[Ink,line width=0.9pt] (2.2*\k,0)--(2.2*\k+3,3);}
\foreach \j in {0,...,3}{\draw[Ink,line width=0.9pt] (\j,\j)--(\j+11,\j);}
\end{scope}\end{figC}
\begin{enumerate}[label=\alph*),leftmargin=1.6em,itemsep=1pt]
\item Trong hình có tất cả bao nhiêu hình bình hành?
\item Có bao nhiêu hình bình hành chứa hình bình hành nhỏ được tô màu?
\item Có bao nhiêu hình bình hành được ghép từ \emph{đúng} $4$ hình bình hành nhỏ?
\end{enumerate}
\end{probbox}
\probsol{
\textbf{Phân tích đề.} Hình có $4$ đường ngang và $6$ đường xiên. Như Bài 27, mỗi hình bình hành ứng với một
cách chọn $2$ đường ngang (cạnh trên, cạnh dưới) và $2$ đường xiên (cạnh trái, cạnh phải).
\begin{recallbox}Số cách chọn $2$ đường từ $k$ đường là $\dfrac{k(k-1)}{2}$.\end{recallbox}

\textbf{a)}
\step{1}{Đếm}
Chọn $2$ trong $4$ đường ngang: $\tfrac{4\cdot 3}{2}=6$ cách; chọn $2$ trong $6$ đường xiên: $\tfrac{6\cdot5}{2}=15$
cách. Số hình bình hành: $6\times 15=90$.

\textbf{b)}
\step{2}{Điều kiện để chứa ô tô màu}
Đánh số các đường ngang từ dưới lên là $0,1,2,3$; ô tô màu nằm giữa đường $1$ và đường $2$. Hình bình hành chứa ô
này khi cạnh dưới là đường $0$ hoặc $1$ ($2$ cách) và cạnh trên là đường $2$ hoặc $3$ ($2$ cách).
Đánh số các đường xiên từ trái sang là $0,1,\dots,5$; ô tô màu nằm giữa đường $2$ và đường $3$. Cạnh trái là đường
$0$, $1$ hoặc $2$ ($3$ cách); cạnh phải là đường $3$, $4$ hoặc $5$ ($3$ cách).
\step{3}{Nhân lại}
$2\times 2\times 3\times 3=36$ hình bình hành.
\begin{proofbox}Vì sao được nhân? Bốn lựa chọn (dưới, trên, trái, phải) độc lập với nhau: chọn thế nào cũng được
đúng một hình bình hành chứa ô tô màu, và hình khác nhau ứng với lựa chọn khác nhau.\end{proofbox}

\textbf{c)}
\step{4}{Các ``kích thước'' gồm đúng $4$ ô}
Một hình bình hành cao $h$ hàng, rộng $w$ cột gồm $h\times w$ ô. Cần $h\times w=4$ với $h\le 3$, $w\le 5$:
\begin{itemize}[leftmargin=1.4em,itemsep=1pt]
\item $h=1$, $w=4$: chọn hàng ($3$ cách), chọn $4$ cột liền nhau trong $5$ cột ($2$ cách): $3\cdot 2=6$;
\item $h=2$, $w=2$: chọn $2$ hàng liền nhau ($2$ cách), chọn $2$ cột liền nhau ($4$ cách): $2\cdot 4=8$;
\item $h=4$, $w=1$: không được vì chỉ có $3$ hàng.
\end{itemize}
Tổng: $6+8=14$.
\begin{warnbox}Đừng quên kiểu $2\times 2$! Và hình cao $4$ hàng thì không có.\end{warnbox}
\verify{Câu a) đếm theo hàng: hình cao $1$ hàng: $3\cdot 15=45$; cao $2$ hàng: $2\cdot 15=30$; cao $3$ hàng: $15$
(mỗi dải có $5+4+3+2+1=15$ cách chọn cột). Tổng $90$. Khớp.}
\ans{a) $90$. \ b) $36$. \ c) $14$.}
}

% ---------------------------------------------------------------- 49
\begin{probbox}{49}{\stars{2}}
\textit{(Mở rộng Bài 28, 29.)} Một hình hộp chữ nhật kích thước $4\times 5\times 6$ được ghép từ $120$ khối lập
phương nhỏ cạnh $1$. Người ta sơn tất cả các mặt ngoài của hình hộp, rồi tách rời các khối nhỏ.
\begin{enumerate}[label=\alph*),leftmargin=1.6em,itemsep=1pt]
\item Có bao nhiêu khối nhỏ được sơn đúng $3$ mặt? đúng $2$ mặt? đúng $1$ mặt? không được sơn mặt nào?
\item Tổng số mặt nhỏ (hình vuông $1\times 1$) được sơn là bao nhiêu? Giải thích vì sao kết quả bằng diện tích
toàn phần của hình hộp.
\item Một hình lập phương $n\times n\times n$ ($n\ge 3$) cũng được sơn và tách như vậy. Tìm $n$ để số khối nhỏ được
sơn đúng $1$ mặt bằng số khối nhỏ không được sơn mặt nào.
\end{enumerate}
\end{probbox}
\probsol{
\textbf{Phân tích đề.} Số mặt được sơn của một khối nhỏ phụ thuộc vào \emph{vị trí} của nó: ở đỉnh, trên cạnh,
trên mặt hay nằm hẳn bên trong. Quan sát mặt trước (kích thước $6\times 4$) để thấy rõ:
\begin{figC}\begin{scope}[x=0.5cm,y=0.5cm]
\foreach \x in {0,...,5}{\foreach \z in {0,...,3}{
  \pgfmathtruncatemacro\ex{(\x==0||\x==5)?1:0}\pgfmathtruncatemacro\ez{(\z==0||\z==3)?1:0}
  \pgfmathtruncatemacro\t{\ex+\ez}
  \ifnum\t=2 \fill[tmOrange!60] (\x,\z) rectangle (\x+1,\z+1);\fi
  \ifnum\t=1 \fill[tmGreen!40] (\x,\z) rectangle (\x+1,\z+1);\fi
  \ifnum\t=0 \fill[tmBlue!30] (\x,\z) rectangle (\x+1,\z+1);\fi}}
\draw[Ink] (0,0) grid[step=1] (6,4);
\foreach \x in {0,...,6}{\draw[Ink] (\x,4)--(\x+2.5,5.75);}
\foreach \d in {0,...,5}{\draw[Ink] (0.5*\d,4+0.35*\d)--(6+0.5*\d,4+0.35*\d);}
\foreach \z in {0,...,4}{\draw[Ink] (6,\z)--(8.5,\z+1.75);}
\foreach \d in {0,...,5}{\draw[Ink] (6+0.5*\d,0.35*\d)--(6+0.5*\d,4+0.35*\d);}
\draw[Ink,line width=1.1pt] (0,0)--(6,0)--(8.5,1.75)--(8.5,5.75)--(2.5,5.75)--(0,4)--cycle (0,4)--(6,4)--(6,0) (6,4)--(8.5,5.75);
\node[below,font=\footnotesize] at (3,0){$6$};\node[left,font=\footnotesize] at (0,2){$4$};\node[right,font=\footnotesize] at (7.3,0.6){$5$};
\begin{scope}[xshift=5.2cm,yshift=0.4cm,font=\footnotesize]
\fill[tmOrange!60] (0,3) rectangle (0.6,3.6);\node[right] at (0.7,3.3){đỉnh: $3$ mặt sơn};
\fill[tmGreen!40] (0,2) rectangle (0.6,2.6);\node[right] at (0.7,2.3){trên cạnh: $2$ mặt};
\fill[tmBlue!30] (0,1) rectangle (0.6,1.6);\node[right] at (0.7,1.3){giữa một mặt: $1$ mặt};
\end{scope}
\end{scope}\end{figC}

\textbf{a)}
\step{1}{Khối ở đỉnh: $3$ mặt sơn}
Hình hộp có $8$ đỉnh, mỗi đỉnh có đúng một khối: $8$ khối.
\step{2}{Khối trên cạnh (trừ hai đầu): $2$ mặt sơn}
Hình hộp có $12$ cạnh: $4$ cạnh dài $4$, $4$ cạnh dài $5$, $4$ cạnh dài $6$. Trên một cạnh dài $L$ có $L$ khối, bỏ
$2$ khối ở hai đầu (đã đếm ở đỉnh) còn $L-2$ khối:
\[ 4\cdot(4-2)+4\cdot(5-2)+4\cdot(6-2)=8+12+16=36. \]
\step{3}{Khối ở giữa một mặt: $1$ mặt sơn}
Trên một mặt $p\times q$, bỏ đi viền ngoài còn $(p-2)(q-2)$ ô. Hình hộp có $2$ mặt $4\times5$, $2$ mặt $4\times 6$,
$2$ mặt $5\times 6$:
\[ 2\cdot(2\cdot 3)+2\cdot(2\cdot 4)+2\cdot(3\cdot 4)=12+16+24=52. \]
\step{4}{Khối nằm hẳn bên trong: không sơn}
Bóc đi ``lớp vỏ'' dày $1$ ở mọi phía, còn lại một hình hộp $(4-2)\times(5-2)\times(6-2)=2\times 3\times 4=24$ khối.
\verify{$8+36+52+24=120$ -- đúng bằng tổng số khối.}

\textbf{b)}
\step{5}{Đếm mặt nhỏ được sơn}
$8\cdot 3+36\cdot 2+52\cdot 1+24\cdot 0=24+72+52=148$.
\begin{recallbox}Diện tích toàn phần hình hộp $a\times b\times c$ là $2(ab+bc+ca)$.\end{recallbox}
$2(4\cdot5+5\cdot6+6\cdot4)=2(20+30+24)=148$. Hai kết quả bằng nhau vì mỗi ô vuông $1\times 1$ trên mặt ngoài
hình hộp chính là đúng một mặt nhỏ được sơn của đúng một khối -- ta chỉ đếm cùng một tập hợp theo hai cách.

\textbf{c)}
\step{6}{Lập phương trình}
Với hình lập phương $n\times n\times n$: mỗi mặt có $(n-2)^2$ ô ở giữa, $6$ mặt cho $6(n-2)^2$ khối sơn đúng $1$
mặt; phần bên trong là hình lập phương $(n-2)^3$ khối. Cần:
\[ (n-2)^3=6(n-2)^2. \]
\step{7}{Giải}
Vì $n\ge 3$ nên $(n-2)^2>0$; chia hai vế cho $(n-2)^2$: $n-2=6$, tức $n=8$.
\begin{warnbox}Điều kiện $n\ge 3$ quan trọng: với $n=2$ thì cả hai số đều bằng $0$ (``bằng nhau'' một cách tầm
thường), và ta không được chia cho $0$.\end{warnbox}
\verify{$n=8$: $6\cdot 6^2=216$ và $6^3=216$. Đúng.}
\ans{a) $8$; $36$; $52$; $24$ khối. \ b) $148$. \ c) $n=8$.}
}

% ---------------------------------------------------------------- 50
\begin{probbox}{50}{\stars{2}}
\textit{(Mở rộng Bài 19, 25.)} Ghép $10$ ô vuông cạnh $1$ thành một hình phẳng sao cho: hai ô kề nhau thì ghép
trùng khít một cạnh, và hình ``liền một khối'' (từ ô bất kì có thể đi đến ô bất kì khác qua các cạnh chung).
Chu vi của hình là tổng độ dài các cạnh ô không bị ô khác che.
\begin{figC}\begin{scope}[x=0.38cm,y=0.38cm]
\foreach \c in {(0,0),(1,0),(2,0),(2,1),(3,1),(4,1),(4,2),(1,1),(4,0),(5,1)}{\fill[tmBlue!30,draw=Ink] \c rectangle ++(1,1);}
\node[font=\footnotesize] at (3,-0.8){một ví dụ};
\end{scope}\end{figC}
\begin{enumerate}[label=\alph*),leftmargin=1.6em,itemsep=1pt]
\item Chứng tỏ chu vi luôn là một số chẵn.
\item Chu vi lớn nhất có thể là bao nhiêu?
\item Chu vi nhỏ nhất có thể là bao nhiêu?
\item Những số nào có thể là chu vi của một hình như vậy?
\end{enumerate}
\end{probbox}
\probsol{
\textbf{Phân tích đề.} Gọi $k$ là số \emph{cạnh chung} (cạnh mà hai ô ghép vào nhau). Ý chính: mỗi cạnh chung
làm ``mất'' hai cạnh ô khỏi chu vi.

\textbf{a)}
\step{1}{Công thức chu vi}
$10$ ô có tất cả $10\cdot 4=40$ cạnh ô. Mỗi cạnh chung là cạnh của hai ô và không nằm trên chu vi; mọi cạnh ô
còn lại nằm trên chu vi. Do đó
\[ C=40-2k. \]
$40$ và $2k$ đều chẵn nên $C$ chẵn.

\textbf{b)}
\step{2}{Hình liền khối có ít nhất $9$ cạnh chung}
\begin{proofbox}Ta dựng lại hình bằng cách đặt từng ô một: bắt đầu từ một ô bất kì, mỗi lần đặt thêm một ô kề với
một ô đã đặt (luôn làm được vì hình liền khối). Mỗi ô thêm vào tạo ra ít nhất $1$ cạnh chung mới. Có $9$ lần đặt
thêm nên $k\ge 9$.\end{proofbox}
\step{3}{Kết luận}
$C=40-2k\le 40-18=22$. Dấu bằng xảy ra với hàng thẳng $1\times 10$: chu vi $2(10+1)=22$. Vậy chu vi lớn nhất là $22$.

\textbf{c)}
\step{4}{So sánh với hình chữ nhật bao quanh}
Gọi hình chữ nhật nhỏ nhất (cạnh theo lưới) chứa cả hình là $a\times b$ ($a$ cột, $b$ hàng).
\begin{proofbox}Mỗi cột của hình chữ nhật bao có ít nhất một ô. Ô cao nhất trong cột có cạnh trên nằm trên chu vi,
ô thấp nhất có cạnh dưới nằm trên chu vi: mỗi cột góp ít nhất $2$ đơn vị cạnh ngang. Tương tự mỗi hàng góp ít nhất
$2$ đơn vị cạnh đứng. Vậy $C\ge 2a+2b=2(a+b)$.\end{proofbox}
\step{5}{Chặn $a+b$}
Hình chữ nhật bao chứa $10$ ô nên $ab\ge 10$. Nếu $a+b\le 6$ thì $ab\le 9$ (các cặp có tổng $6$: $1\cdot5=5$,
$2\cdot4=8$, $3\cdot3=9$; tổng nhỏ hơn còn cho tích nhỏ hơn) -- vô lí. Vậy $a+b\ge 7$ và $C\ge 14$.
\step{6}{Có hình đạt $14$}
Hình chữ nhật $2\times 5$: chu vi $2(2+5)=14$. Vậy chu vi nhỏ nhất là $14$.

\textbf{d)}
\step{7}{Các giá trị có thể}
Từ a), b), c): $C$ chẵn và $14\le C\le 22$, nên $C\in\{14;16;18;20;22\}$. Mỗi giá trị đều đạt được bằng các hình
``hai hàng'' sau (dùng ý của Bài 25: đường viền là cầu thang nên chu vi bằng chu vi hình chữ nhật bao quanh):
\begin{figC}\begin{scope}[x=0.3cm,y=0.3cm]
\foreach \top/\bot/\xs/\lab in {5/5/0/14,6/4/6.5/16,7/3/14/18,8/2/22.5/20}{
  \foreach \i in {1,...,\top}{\fill[tmBlue!30,draw=Ink] (\xs+\i-1,1) rectangle ++(1,1);}
  \foreach \i in {1,...,\bot}{\fill[tmBlue!30,draw=Ink] (\xs+\i-1,0) rectangle ++(1,1);}
  \node[font=\footnotesize] at (\xs+\top/2,-0.9){$C=\lab$};}
\foreach \i in {1,...,10}{\fill[tmGreen!35,draw=Ink] (\i-1,-4) rectangle ++(1,1);}
\node[font=\footnotesize,anchor=west] at (10.5,-3.5){$C=22$};
\end{scope}\end{figC}
Hàng trên $5,6,7,8$ ô, hàng dưới $5,4,3,2$ ô: hình chữ nhật bao là $5\times2$, $6\times 2$, $7\times 2$, $8\times 2$,
chu vi lần lượt $14,16,18,20$; và hàng thẳng cho $22$.
\begin{warnbox}Không phải cứ ``gọn'' là chu vi nhỏ: hình $3\times 3$ thêm một ô lồi ra có chu vi $12+2=14$, bằng
hình $2\times 5$, chứ không nhỏ hơn.\end{warnbox}
\verify{Hình $6+4$: $k=5+3+4=12$, $C=40-24=16$. Hình $8+2$: $k=7+1+2=10$, $C=40-20=20$. Đúng.}
\ans{a) $C=40-2k$ chẵn. \ b) $22$. \ c) $14$. \ d) $14$, $16$, $18$, $20$, $22$.}
}
